% Options for packages loaded elsewhere
\PassOptionsToPackage{unicode}{hyperref}
\PassOptionsToPackage{hyphens}{url}
\PassOptionsToPackage{dvipsnames,svgnames,x11names}{xcolor}
\documentclass[
  10pt,
  fontset=fandol]{ctexart}
\usepackage{xcolor}
\usepackage[margin=16mm]{geometry}
\usepackage{amsmath,amssymb}
\setcounter{secnumdepth}{-\maxdimen} % remove section numbering
\usepackage{iftex}
\ifPDFTeX
  \usepackage[T1]{fontenc}
  \usepackage[utf8]{inputenc}
  \usepackage{textcomp} % provide euro and other symbols
\else % if luatex or xetex
  \usepackage{unicode-math} % this also loads fontspec
  \defaultfontfeatures{Scale=MatchLowercase}
  \defaultfontfeatures[\rmfamily]{Ligatures=TeX,Scale=1}
\fi
\usepackage{lmodern}
\ifPDFTeX\else
  % xetex/luatex font selection
\fi
% Use upquote if available, for straight quotes in verbatim environments
\IfFileExists{upquote.sty}{\usepackage{upquote}}{}
\IfFileExists{microtype.sty}{% use microtype if available
  \usepackage[]{microtype}
  \UseMicrotypeSet[protrusion]{basicmath} % disable protrusion for tt fonts
}{}
\makeatletter
\@ifundefined{KOMAClassName}{% if non-KOMA class
  \IfFileExists{parskip.sty}{%
    \usepackage{parskip}
  }{% else
    \setlength{\parindent}{0pt}
    \setlength{\parskip}{6pt plus 2pt minus 1pt}}
}{% if KOMA class
  \KOMAoptions{parskip=half}}
\makeatother
\setlength{\emergencystretch}{3em} % prevent overfull lines
\providecommand{\tightlist}{%
  \setlength{\itemsep}{0pt}\setlength{\parskip}{0pt}}
\usepackage{amsmath,amssymb,mathtools,needspace}
\xeCJKDeclareCharClass{CJK}{"2032,"0400->"04FF,"2160->"216B,"2460->"2473,"25A0,"25C7,"25CB,"25CF}
\IfFontExistsTF{Arial Unicode MS}{\xeCJKsetup{AutoFallBack=true}\setCJKfallbackfamilyfont{\CJKrmdefault}{Arial Unicode MS}}{}
\setcounter{secnumdepth}{0}
\setlength{\emergencystretch}{2em}
\allowdisplaybreaks
\usepackage{bookmark}
\IfFileExists{xurl.sty}{\usepackage{xurl}}{} % add URL line breaks if available
\urlstyle{same}
\hypersetup{
  pdftitle={矩阵的三角分解},
  colorlinks=true,
  linkcolor={Maroon},
  filecolor={Maroon},
  citecolor={Blue},
  urlcolor={Blue},
  pdfcreator={LaTeX via pandoc}}

\title{矩阵的三角分解}
\author{}
\date{}

\begin{document}
\maketitle

\subsubsection{7.2.2 矩阵的三角分解}\label{ux77e9ux9635ux7684ux4e09ux89d2ux5206ux89e3}

下面借助矩阵理论进一步对消去法作些分析，从而建立 Gauss 消去法与矩阵因式分解的关系。

设式（7.2.1）中 \(A\) 的各顺序主子式均不为零。由于对 \(A\) 施行行的初等变换相当于用初等矩阵左乘 \(A\)，于是对式（7.2.1）施行第一次消元后化为式（7.2.7），这时 \(A^{(1)}\) 化为 \(A^{(2)}\)，\(b^{(1)}\) 化为 \(b^{(2)}\)，即

\[L_1A^{(1)}=A^{(2)},\qquad L_1b^{(1)}=b^{(2)},\]

其中

\[
L_1=\begin{pmatrix}
1&&&&\\
-m_{21}&1&&&\\
-m_{31}&&1&&\\
\vdots&&&\ddots&\\
-m_{n1}&&&&1
\end{pmatrix}.
\]

\href{../../source-images/pdf-182.jpeg}{原页扫描}

一般第 \(k\) 步消元，\(A^{(k)}\) 化为 \(A^{(k+1)}\)，\(b^{(k)}\) 化为 \(b^{(k+1)}\)，相当于

\[L_kA^{(k)}=A^{(k+1)},\qquad L_kb^{(k)}=b^{(k+1)}.\]

重复这一过程，最后得到

\[
\begin{cases}
L_{n-1}\cdots L_2L_1A^{(1)}=A^{(n)},\\
L_{n-1}\cdots L_2L_1b^{(1)}=b^{(n)},
\end{cases}
\tag{7.2.14}
\]

其中

\[
L_k=\begin{pmatrix}
1&&&&&\\
&\ddots&&&&\\
&&1&&&\\
&&-m_{k+1,k}&1&&\\
&&\vdots&&\ddots&\\
&&-m_{nk}&&&1
\end{pmatrix}.
\]

将上三角矩阵 \(A^{(n)}\) 记作 \(U\)，由式（7.2.14）得到

\[A=L_1^{-1}L_2^{-1}\cdots L_{n-1}^{-1}U=LU,\]

其中

\[
L=L_1^{-1}L_2^{-1}\cdots L_{n-1}^{-1}
=\begin{pmatrix}
1&&&&\\
m_{21}&1&&&\\
m_{31}&m_{32}&1&&\\
\vdots&\vdots&\ddots&\ddots&\\
m_{n1}&m_{n2}&\cdots&m_{n,n-1}&1
\end{pmatrix}.
\]

为单位下三角矩阵。

这就是说，Gauss 消去法实质上产生了一个将 \(A\) 分解为两个三角矩阵相乘的因式分解，于是得到如下重要定理，它在解方程组的直接法中起着重要作用。

\textbf{定理 7.3（矩阵的 LU 分解）} 设 \(A\) 为 \(n\) 阶矩阵，如果 \(A\) 的顺序主子式 \(D_i\ne0\)（\(i=1,2,\cdots,n-1\)），则 \(A\) 可分解为一个单位下三角矩阵 \(L\) 和一个上三角矩阵 \(U\) 的乘积，且这种分解是唯一的。

\textbf{证明} 根据以上 Gauss 消去法的矩阵分析，\(A=LU\) 的存在性已经得到证明，现仅在 \(A\) 为非奇异矩阵的假定下来证明它的唯一性，当 \(A\) 为奇异矩阵的情况留作课外练习。

设

\[A=LU=L_1U_1,\]

其中 \(L,L_1\) 为单位下三角矩阵；\(U,U_1\) 为上三角矩阵。

由于 \(U_1^{-1}\) 存在，故

\[L^{-1}L_1=UU_1^{-1}.\]

上式右端为上三角矩阵，左端为单位下三角矩阵，从而上式两端都必须等于单位矩阵，故 \(U=U_1\)，\(L=L_1\)。证毕。

\textbf{例 7.2} 对于例 7.1，系数矩阵 \(A=\begin{pmatrix}1&1&1\\0&4&-1\\2&-2&1\end{pmatrix}\)，由 Gauss 消去法，有

\href{../../source-images/pdf-183.jpeg}{原页扫描}

\[m_{21}=0,\quad m_{31}=2,\quad m_{32}=-1,\]

故

\[
A=\begin{pmatrix}1&0&0\\0&1&0\\2&-1&1\end{pmatrix}
\begin{pmatrix}1&1&1\\0&4&-1\\0&0&-2\end{pmatrix}=LU.
\]

\end{document}
